%    Latex Template for ACA 2016
%%%%%%%%%%%%%%%%%%%%%%%%%%%%%%%%%%%%%%%%
%    DO NOT CHANGE THE SETTINGS BELOW
%%%%%%%%%%%%%%%%%%%%%%%%%%%%%%%%%%%%%%%%
\documentclass[11pt]{article}
\usepackage{epsfig}
\usepackage{graphicx}
\usepackage[latin1]{inputenc}
\usepackage[T1]{fontenc}
\usepackage{mathptmx}
\usepackage{url}



\begin{document}



%%%%%%%%%%%%%%%%%%%%%%%%
%%%%%%%%%%%%%%%%%%%%%%%%%%%%%%%%%%%%%%%%
%    DO NOT CHANGE THE SETTINGS ABOVE
%%%%%%%%%%%%%%%%%%%%%%%%%%%%%%%%%%%%%%%%
%
\vspace{-0.8cm}
\begin{flushleft}
\Large \textbf{\noindent
%%% PLEASE INSERT THE TITLE OF YOUR TALK HERE %%%%%%%
Applications of  CAS  in the Teaching and Learning of Discrete Mathematics}
%%%%%%%%%%%%%%%%%%%%%%%%%%%%%%%%%%%%%%%%%%%%%%%%%%%%%%%%
\\
%%%%%%%%%%%%%%%%%%%%%%%%%%%%%%
\vspace{0.5cm}
\normalsize
 %%%%%%             NAMES OF AUTHORS         %%%%%%
 %%%%%% NAME OF SPEAKER SHOULD BE UNDERLINED %%%%%%%%%
\normalsize{
 %%%%%% FOR EXAMPLE %%%%%
 \underline{M. Durcheva}$^1$, E. Varbanova$^2$
 %%%%%%%%%%%%%%%%%%%%%%%%
} \\
\vspace{5mm}
\textit{\footnotesize
 %%%%%% AFFILIATION OF AUTHORS %%%%%%
$^1$ Technical University of Sofia, Bulgaria,
mdurcheva66@gmail.com\\
$^2$ Technical University of Sofia, Bulgaria,
elvar@tu-sofia.bg\\
 %%%%%%%%%%%%%%%%%%%%%%%%%%%%%%%%%%%%%
}
\end{flushleft}

      Discrete Mathematics is a compulsory course in the curriculum of the Bachelor Degree Program at the Faculty of Applied Mathematics and Informatics of the Technical University of Sofia (TUS). Discrete structures are taught also to engineering students at TUS.

As it is reported in  Refs.~\cite{FRAM,NISS}), the following competences of the above mentioned students are to be developed: thinking mathematically; reasoning mathematically; posing and solving mathematical problems; modeling mathematically; representing mathematical entities; handling mathematical symbols and formalism; communicating in, with, and about mathematics; making use of aids and tools. How could applications of CAS contribute to this educational goal?
It is often necessary to manipulate big data in discrete mathematics, mainly for the purpose of combinatorics, recurrence relations, formal logic, Boolean algebra. For that reason formulas have very limited range of use.  The application of CAS can help expand these limits significantly. The most important thing is how learning can be enhanced by the use of CAS.

 Concurrent implementation of pedagogy and technology requires consideration of current practices, including the activities and expectations of learners and teachers. Technology allows students to work with more complex and realistic models since work within the model is supported by technology and even the setup of models can be facilitated by technology. In engineering programs, models are often only partially visible, so students need to learn to work with technology where the underlying model is not known to them. This requires knowledge about strategies for checking one's understanding of the workings of the program and also for checking the results.

In this paper our experience in application of CAS MAPLE for modeling discrete structures is represented. Some procedures for calculating combinations with repetition and variations with repetition in case of arbitrary restrictions on appearance of elements are considered. The result is the list of the allowed number of appearances of the element in the polynomial involved in the generated function  (see Refs.~\cite{GEN}). Some conclusions are visualized. Solving recurrence equations with MAPLE ( Refs.~\cite{ROS})) is also discussed.








%%%%%%%%%%%%%%%%%%%%%%%%%%%%%%%%%%%%%%%


%%%%%%%%%%%%%%%%%%%%%%%%%%%%%%%%%%%%%%%%
%% BIBLIOGRAPHY
%%%%%%%%%%%%%%%%%%%%%%%%%%%%%%%%%%%%%%%%
\begin{thebibliography}{00}
\addcontentsline{toc}{chapter}{References}
\setlength{\itemsep}{-1mm}
\footnotesize
\bibitem{FRAM} 	\textit{A Framework for Mathematics Curricula in Engineering Education}, A Report of the Mathematics Working Group, European Society for Engineering Education (SEFI), Brussels, 2013, Principal Editor: Burkhard Alpers.
\bibitem{NISS} 	M. Niss and T. Hojgaard, \textit{Competencies and Mathematical Learning Ideas and inspiration for the development of mathematics teaching and learning in Denmark}, IMFUFA tekst nr. 485/ 2011.
\bibitem{ROS} K. Rosen, \textit{Discrete Mathematics and Its Applications}, Seventh ed., 2016, McGraw-Hill Global Education Holdings, LLC.
\bibitem{GEN} H. Wilf, \textit{Generatingfunctionology}, Academic Press, Inc., New York, 1994.
\end{thebibliography}
\normalsize
\end{document}
